% !TEX program = xelatex
\documentclass[a4paper,11pt]{article}

% ===================== Packages =====================
\usepackage{fontspec}
\usepackage{xeCJK}
\usepackage[a4paper,top=2.8cm,bottom=2.8cm,left=2.6cm,right=2.6cm,headheight=15pt,headsep=10pt,footskip=22pt]{geometry}
\usepackage{array}
\usepackage{amsmath}
\usepackage{amssymb}
\usepackage{booktabs}
\usepackage{tabularx}
\usepackage{longtable}
\usepackage{enumitem}
\usepackage{xcolor}
\usepackage{setspace}
\usepackage{titlesec}
\usepackage{titletoc}
\usepackage{float}
\usepackage{tikz}
\usepackage{fancyhdr}
\usepackage[font=small,labelfont=bf,labelsep=quad,skip=6pt]{caption}
\usepackage{listings}
\usepackage{hyperref}

\usetikzlibrary{arrows.meta,positioning,fit,calc,shapes.geometric,backgrounds}

% ===================== Fonts =====================
\IfFontExistsTF{Noto Sans CJK SC}{%
  \setCJKmainfont{Noto Sans CJK SC}%
  \setCJKsansfont{Noto Sans CJK SC}%
  \setCJKmonofont{Noto Sans CJK SC}%
}{%
  \IfFontExistsTF{PingFang SC}{%
    \setCJKmainfont{PingFang SC}%
    \setCJKsansfont{PingFang SC}%
    \setCJKmonofont{PingFang SC}%
  }{%
    \setCJKmainfont{AR PL UMing CN}%
    \setCJKsansfont{AR PL UKai CN}%
    \setCJKmonofont{AR PL UMing CN}%
  }%
}
\IfFontExistsTF{Noto Serif CJK SC}{%
  \setmainfont{Noto Serif CJK SC}%
}{%
  \IfFontExistsTF{Songti SC}{\setmainfont{Songti SC}}{\setmainfont{DejaVu Serif}}%
}
\IfFontExistsTF{Noto Sans CJK SC}{%
  \setsansfont{Noto Sans CJK SC}%
}{%
  \IfFontExistsTF{PingFang SC}{\setsansfont{PingFang SC}}{\setsansfont{DejaVu Sans}}%
}
\IfFontExistsTF{Noto Sans Mono CJK SC}{%
  \setmonofont{Noto Sans Mono CJK SC}%
}{%
  \IfFontExistsTF{Menlo}{\setmonofont{Menlo}}{\setmonofont{DejaVu Sans Mono}}%
}

% \XeTeXlinebreaklocale "zh"  % Tectonic's bundled ICU may not provide this locale
\XeTeXlinebreakskip = 0pt plus 1pt

% ===================== Hyperref =====================
\hypersetup{
  colorlinks=true,
  linkcolor=black,
  urlcolor=black,
  citecolor=black,
  pdftitle={ForL0 State Backend 使用说明书 v2.0},
  pdfauthor={ForL0 Project}
}

% ===================== Localization =====================
\renewcommand{\contentsname}{目\quad 录}
\renewcommand{\tablename}{表}
\renewcommand{\figurename}{图}

% ===================== Spacing =====================
\setstretch{1.28}
\setlength{\parindent}{2em}
\setlength{\parskip}{0.35em}
\setlength{\emergencystretch}{3em}
\sloppy
\setlist[itemize]{leftmargin=2em,itemsep=0.25em,topsep=0.35em,parsep=0pt}
\setlist[enumerate]{leftmargin=2.2em,itemsep=0.25em,topsep=0.35em,parsep=0pt}

% ===================== Headings =====================
\definecolor{accent}{RGB}{30,64,110}
\definecolor{rulegray}{RGB}{60,60,60}
\definecolor{softgray}{RGB}{245,246,248}
\definecolor{midgray}{RGB}{220,223,228}
\definecolor{linegray}{RGB}{120,124,132}
\definecolor{codebg}{RGB}{248,249,251}
\definecolor{codeborder}{RGB}{220,223,228}
\definecolor{codekw}{RGB}{30,64,110}
\definecolor{codecomment}{RGB}{90,100,110}
\definecolor{codestring}{RGB}{102,70,20}

\titleformat{\section}[hang]
  {\normalfont\large\bfseries\color{accent}}
  {\thesection}{0.8em}{}
  [\vspace{2pt}{\color{accent!60}\titlerule[0.6pt]}]
\titlespacing*{\section}{0pt}{1.6em}{0.9em}

\titleformat{\subsection}[hang]
  {\normalfont\normalsize\bfseries\color{accent!85!black}}
  {\thesubsection}{0.6em}{}
\titlespacing*{\subsection}{0pt}{1.2em}{0.55em}

\titleformat{\subsubsection}[hang]
  {\normalfont\normalsize\bfseries}
  {\thesubsubsection}{0.55em}{}
\titlespacing*{\subsubsection}{0pt}{0.9em}{0.4em}

% ===================== TOC styling =====================
\titlecontents{section}
  [0em]{\addvspace{0.55em}\bfseries}
  {\contentslabel{1.6em}}{}
  {\titlerule*[0.6pc]{.}\contentspage}
\titlecontents{subsection}
  [2.2em]{\addvspace{0.15em}}
  {\contentslabel{2.4em}}{}
  {\titlerule*[0.6pc]{.}\contentspage}
\titlecontents{subsubsection}
  [4.6em]{\addvspace{0.05em}\small}
  {\contentslabel{3.0em}}{}
  {\titlerule*[0.6pc]{.}\contentspage}

% ===================== Header / Footer =====================
\newcommand{\project}{ForL0 State Backend}
\pagestyle{fancy}
\fancyhf{}
\renewcommand{\headrulewidth}{0.4pt}
\renewcommand{\footrulewidth}{0pt}
\fancyhead[L]{\small\sffamily\color{accent}\project{} 使用说明书}
\fancyhead[R]{\small\sffamily\color{accent}v2.0}
\fancyfoot[C]{\small\thepage}

\fancypagestyle{plain}{%
  \fancyhf{}%
  \renewcommand{\headrulewidth}{0pt}%
  \fancyfoot[C]{\small\thepage}%
}

% ===================== Listings =====================
\lstdefinestyle{codestyle}{
  basicstyle=\ttfamily\small,
  backgroundcolor=\color{codebg},
  frame=single,
  rulecolor=\color{codeborder},
  framesep=4pt,
  xleftmargin=6pt,
  xrightmargin=4pt,
  aboveskip=0.6em,
  belowskip=0.6em,
  breaklines=true,
  breakatwhitespace=false,
  showstringspaces=false,
  columns=fullflexible,
  keepspaces=true,
  tabsize=2,
  commentstyle=\color{codecomment}\itshape,
  stringstyle=\color{codestring},
  keywordstyle=\color{codekw}\bfseries,
  extendedchars=false,
  inputencoding=utf8,
  literate={\$}{{\$}}1 {\&}{{\&}}1 {\_}{{\_}}1
}
\lstdefinelanguage{shellcmd}{
  morekeywords={cd,ls,mkdir,cp,mv,rm,export,cmake,make,mvn,yum,sudo,wget,tar,ln,grep,nm,file,ldd,source,ln,tee,curl,echo,python3,pip3,bash,nproc},
  morecomment=[l]{\#},
  morestring=[b]",
}
\lstdefinelanguage{yamlcmd}{
  morekeywords={state,backend,type,forl0,checkpoints,dir,env,java,opts,all,l0-cache,enabled,size,max-per-alloc,async-snapshots,execution,checkpointing,interval,mode,savepoints,checkpoint-storage,externalized-checkpoint-retention},
  morecomment=[l]{\#},
  morestring=[b]",
}
\lstdefinelanguage{javalang}{
  morekeywords={public,private,protected,class,void,new,return,this,if,else,for,while,final,static,import,package,extends,implements,throws,try,catch,throw,null,true,false,long,int,boolean,String,Override},
  morecomment=[l]{//},
  morecomment=[s]{/*}{*/},
  morestring=[b]",
}
\lstset{style=codestyle,language=shellcmd}

% ===================== Table column types =====================
\newcolumntype{L}[1]{>{\raggedright\arraybackslash}p{#1}}
\newcolumntype{C}[1]{>{\centering\arraybackslash}p{#1}}
\newcommand{\cfg}[1]{\nolinkurl{#1}}

% ===================== Figure palette =====================
\definecolor{fgfill}{RGB}{248,249,251}
\definecolor{fgaccent}{RGB}{30,64,110}
\definecolor{fgaccentlight}{RGB}{223,231,242}
\definecolor{fgbox}{RGB}{255,255,255}
\definecolor{fgborder}{RGB}{90,95,105}
\definecolor{fgsubtle}{RGB}{235,237,240}
\definecolor{fgline}{RGB}{70,74,82}

\tikzset{
  figfont/.style={font=\footnotesize\sffamily},
  figsmall/.style={font=\scriptsize\sffamily,text=black!70},
  figtitle/.style={font=\footnotesize\sffamily\bfseries,text=fgaccent},
  box/.style={draw=fgborder, line width=0.5pt, rounded corners=2pt, fill=fgbox,
              minimum height=0.95cm, inner xsep=10pt, inner ysep=5pt, align=center, figfont},
  accentbox/.style={draw=fgaccent, line width=0.6pt, rounded corners=2pt, fill=fgaccentlight,
              minimum height=0.95cm, inner xsep=10pt, inner ysep=5pt, align=center, figfont},
  subtlebox/.style={draw=fgborder, line width=0.5pt, rounded corners=2pt, fill=fgsubtle,
              minimum height=0.95cm, inner xsep=10pt, inner ysep=5pt, align=center, figfont},
  lane/.style={draw=fgborder!55, line width=0.4pt, rounded corners=3pt, fill=fgfill},
  line/.style={-{Latex[length=2mm,width=1.6mm]}, draw=fgline, line width=0.55pt},
  thinline/.style={draw=fgline!80, line width=0.4pt}
}

% ===================== Title page macro =====================
\newcommand{\coverrule}{{\color{accent}\rule{\linewidth}{0.9pt}}}
\newcommand{\coverrulethin}{{\color{accent!55}\rule{\linewidth}{0.3pt}}}

\begin{document}

% =================================================
%                 COVER
% =================================================
\begin{titlepage}
\thispagestyle{empty}
\centering
\vspace*{5.5cm}

\coverrule
\vspace{2pt}
\coverrulethin

\vspace{1.6cm}
{\sffamily\huge\bfseries\color{accent} \project\par}
\vspace{0.8cm}
{\sffamily\LARGE\bfseries 使用说明书\par}

\vspace{1.6cm}
\coverrulethin
\vspace{2pt}
\coverrule

\vfill

{\sffamily\large\color{accent!85} 2026 年 8 月\par}

\vspace{1.2cm}
\end{titlepage}

% =================================================
%                 TOC
% =================================================
\pagenumbering{roman}
\pagestyle{plain}
\tableofcontents
\clearpage
\pagenumbering{arabic}
\setcounter{page}{1}
\pagestyle{fancy}

% =================================================
%                 CONTENT
% =================================================
\section{产品概述}

\project{} 是面向 Apache Flink 的高性能 Keyed State 后端，针对鲲鹏服务器的体系结构特性进行了定向优化。系统以 C++ 原生引擎承担状态的权威存储与访问，以 Java 控制层与 JNI 桥接层承担 Flink 接口适配与状态注册，并通过 SwissTable 哈希结构、SWAR 并行匹配以及可选的 L0 Cache 加速缩短热点状态访问路径。

本文档内容限定在状态后端的配置、作业集成、运行监控与调优。底层架构细节由《设计说明书》给出，编译与部署由《编译构建指导书》给出。

\subsection{关键特性}

\begin{itemize}
  \item 支持本项目已验证的 Flink Keyed State 接口与语义，现有作业代码无需改动。
  \item 原生 SwissTable 哈希结构 + NEON SWAR 并行匹配，针对鲲鹏 ARM64 做指令级优化。
  \item 可选启用鲲鹏 L0 Cache 作为热键访问的上层快路径，硬件条件不满足时自动保持关闭，底层 SwissTable 直接接管状态访问。
  \item 支持 Flink Checkpoint 与 Savepoint，采用完整 KeyGroup 快照和恢复。
  \item 原生 Copy-On-Write 一致性视图，快照数据在异步阶段写出。
\end{itemize}

\subsection{适用范围}

\begin{table}[H]
\centering
\small
\begin{tabularx}{\textwidth}{L{3.2cm} X}
\toprule
\textbf{项目} & \textbf{要求} \\
\midrule
处理器 & 鲲鹏系列处理器 \\
操作系统 & openEuler \\
JDK & JDK 1.8 及以上（aarch64），建议使用华为毕昇 JDK \\
Flink & Apache Flink 1.20.3 \\
运行模式 & 标准模式 / L0 HotCache 模式（需 \texttt{/dev/hisi\_l0} 或 \texttt{/dev/l0}，以及 \texttt{libl0mempool.so}） \\
\bottomrule
\end{tabularx}
\end{table}

\subsection{支持的状态类型}

\project{} 对以下五类常用 Keyed State 提供原生存储实现：

\begin{table}[H]
\centering
\small
\begin{tabularx}{\textwidth}{L{3.2cm} X L{4.6cm}}
\toprule
\textbf{状态类型} & \textbf{用途} & \textbf{实现类} \\
\midrule
ValueState & 单值键控状态 & \texttt{ForL0ValueState} \\
ListState & 列表键控状态 & \texttt{ForL0ListState} \\
MapState & 键控内部映射状态 & \texttt{ForL0MapState} \\
ReducingState & 归约键控状态 & \texttt{ForL0ReducingState} \\
AggregatingState & 聚合键控状态 & \texttt{ForL0AggregatingState} \\
\bottomrule
\end{tabularx}
\end{table}

% ---------- 图：作业与后端的集成关系 ----------
\begin{figure}[htbp]
\centering
\resizebox{\textwidth}{!}{%
\begin{tikzpicture}[x=1cm,y=1cm]
  \node[lane, minimum width=14.8cm, minimum height=4.8cm] at (8.4,2.3) {};

  \node[figtitle, anchor=east] at (1.4,4.1) {用户作业};
  \node[figtitle, anchor=east] at (1.4,2.3) {状态后端};
  \node[figtitle, anchor=east] at (1.4,0.7) {原生数据面};

  \node[accentbox, text width=11.6cm] (job) at (8.4,4.1) {DataStream / SQL 作业（ValueState, ListState, MapState, \ldots）};

  \node[box, text width=3.2cm] (backend) at (3.3,2.3) {ForL0StateBackend};
  \node[box, text width=3.2cm] (keyed)   at (8.4,2.3) {Keyed State Backend};
  \node[box, text width=3.2cm] (ckpt)    at (13.5,2.3) {Checkpoint\,/\\Savepoint};

  \node[accentbox, text width=11.6cm] (native) at (8.4,0.7) {NativeEngine + StateTable + SwissTable + L0 HotCache};

  \draw[line] (job.south)     -- (keyed.north);
  \draw[line] (backend.east)  -- (keyed.west);
  \draw[line] (keyed.east)    -- (ckpt.west);
  \draw[line] (keyed.south)   -- (native.north);
\end{tikzpicture}%
}
\caption{用户作业与 \project{} 的运行时集成关系}
\label{fig:usage-overview}
\end{figure}

\section{启用与部署}

本节假设 \project{} 的 JAR 与原生库已按《编译构建指导书》完成构建。启用流程主要包括将产物放置到 Flink 对应目录、追加配置并重启集群。

\subsection{产物部署}

\begin{lstlisting}[language=shellcmd]
# 1. 状态后端 JAR
cp flink-statebackend-forL0-1.0-SNAPSHOT.jar $FLINK_HOME/lib/

# 2. 原生引擎共享库
mkdir -p $FLINK_HOME/native
cp libforl0_engine.so $FLINK_HOME/native/

# 3. 多节点集群下，将上述两份文件同步到所有 TaskManager 节点
\end{lstlisting}

\subsection{集群配置}

在 \texttt{\$FLINK\_HOME/conf/config.yaml} 中追加如下配置项：

\begin{lstlisting}[language=yamlcmd]
# 启用 ForL0 作为默认状态后端（使用 StateBackendFactory 的全限定类名）
state:
  backend:
    type: org.apache.flink.state.forl0.ForL0StateBackendFactory
  checkpoints:
    dir: file:///data/flink/checkpoints
  checkpoint-storage: filesystem

# 指定原生库路径（针对 TaskManager 进程）
env:
  java:
    opts:
      all: -Djava.library.path=/opt/flink-1.20.3/native
\end{lstlisting}

配置完成后重启集群：

\begin{lstlisting}[language=shellcmd]
$FLINK_HOME/bin/stop-cluster.sh
$FLINK_HOME/bin/start-cluster.sh
\end{lstlisting}

\subsection{加载校验}

集群启动后，可通过 TaskManager 日志确认状态后端与原生引擎是否正确加载：

\begin{lstlisting}[language=shellcmd]
grep "\[ForL0\]" $FLINK_HOME/log/flink-*-taskexecutor-*.log
\end{lstlisting}

输出中至少应包含原生库加载成功一条记录；如已启用 \texttt{l0-cache.enabled=true}，还会在 KeyedStateBackend 创建时输出一条包含 \texttt{l0Enabled=true/false} 的引擎创建记录：

\begin{lstlisting}[language=shellcmd]
[ForL0] Native engine loaded via java.library.path
[ForL0] C++ engine created: handle=..., keyGroups=[0, 8), total=128, l0Enabled=true
\end{lstlisting}

若日志中的 \texttt{l0Enabled=false}，则表示运行环境不满足 L0 硬件门禁或配置未启用，HotCache 快路径未激活，状态访问仍由 SwissTable 正常服务。

\section{配置参考}

\project{} 的配置分为两类：Flink 标准状态后端配置和 \project{} 自身引入的 \texttt{state.backend.forl0.*} 专用配置。两类配置都可在集群配置文件中声明，也可在作业提交时通过 \texttt{-D} 参数覆盖。

\subsection{启用后端}

\subsubsection{配置文件方式（推荐）}

在集群配置文件中声明默认状态后端后，全集群作业自动生效。\textbf{\texttt{state.backend.type} 必须填写 \texttt{StateBackendFactory} 的全限定类名}：

\begin{lstlisting}[language=yamlcmd]
state:
  backend:
    type: org.apache.flink.state.forl0.ForL0StateBackendFactory
\end{lstlisting}

\subsubsection{代码方式}

在作业中显式设置，仅对当前作业生效：

\begin{lstlisting}[language=javalang]
StreamExecutionEnvironment env =
    StreamExecutionEnvironment.getExecutionEnvironment();
env.setStateBackend(new ForL0StateBackend());
\end{lstlisting}

若需通过 Factory 由配置驱动创建，可使用：

\begin{lstlisting}[language=javalang]
env.setStateBackend(
    new ForL0StateBackendFactory().createFromConfig(
        new Configuration(),
        Thread.currentThread().getContextClassLoader()));
\end{lstlisting}

\subsection{配置项列表}

\begin{table}[H]
\centering
\footnotesize
\begin{tabularx}{\textwidth}{L{6.2cm} L{1.6cm} L{1.6cm} X}
\toprule
\textbf{配置项} & \textbf{类型} & \textbf{默认值} & \textbf{说明} \\
\midrule
\cfg{state.backend.type} & String & -- & 填写 \texttt{StateBackendFactory} 的全限定类名以启用本后端，具体值见 \S 3.1。 \\
\cfg{state.backend.forl0.async-snapshots} & boolean & \texttt{true} & 使用原生 Copy-On-Write 异步快照。 \\
\cfg{state.backend.forl0.l0-cache.enabled} & boolean & \texttt{false} & 是否启用 L0 Cache 加速，需鲲鹏硬件支持。 \\
\cfg{state.backend.forl0.l0-cache.size} & MemorySize & 20\,mb & 兼容配置项，表示单个 TaskManager 进程的 L0 配额。 \\
\cfg{state.backend.forl0.l0-cache.total-size} & MemorySize & -- & 主机/设备级 L0 总预算；按 \texttt{expected-engines} 划分给 TaskManager 进程。 \\
\cfg{state.backend.forl0.l0-cache.expected-engines} & int & 1 & 预期进程级 manager 数；应填写 TaskManager 进程数，不是 slot 或 StateEngine 数。 \\
\cfg{state.backend.forl0.l0-cache.strict-allocation} & boolean & \texttt{false} & 为 true 时，硬件门禁或分配失败直接使作业失败，适用于要求 L0 必须生效的部署。 \\
\cfg{state.backend.forl0.l0-cache.state-size} & MemorySize & 1\,mb & 每个符合条件的标量 ValueState 请求的缓存配额。 \\
\cfg{state.backend.forl0.l0-cache.write-bypass-threshold} & long & 1,048,576 & 连续写入达到该次数且无读取时暂停缓存准入；0 表示禁用该策略。 \\
\cfg{state.backend.forl0.native-memory.max-size} & MemorySize & 0 & 单 StateEngine 的 SwissTable native 内存上限；0 表示不设限。 \\
\cfg{forL0.metricsCollector.enabled} & boolean & \texttt{false} & 是否注册 HotCache gauge；需要监控缓存状态时显式开启。 \\
\cfg{state.backend.forl0.initial-table-capacity} & int & 64 & 单张 SwissTable 的初始 slot 数，必须为 2 的幂且 $\ge 8$。 \\
\cfg{state.backend.forl0.max-table-capacity} & int & 0 & 单张 SwissTable 的容量上限；0 表示不设上限。 \\
\cfg{state.backend.forl0.main-table.load-factor-threshold} & double & 0.875 & SwissTable 扩容阈值，取值范围为 0.5--0.875。 \\
\bottomrule
\end{tabularx}
\end{table}

\subsection{L0 HotCache 启用条件}

当 \texttt{l0-cache.enabled} 配置为 \texttt{true} 时，原生引擎在运行期继续检查以下三项：

\begin{enumerate}
  \item 共享库 \texttt{libl0mempool.so} 可通过 \texttt{dlopen} 加载（依次检查 \texttt{libl0mempool.so}、\texttt{/usr/lib}、\texttt{/usr/lib64}、\texttt{/lib}、\texttt{/lib64}）。
  \item L0 设备节点 \texttt{/dev/hisi\_l0} 或 \texttt{/dev/l0} 存在且可访问。
  \item 首次探测性 \texttt{l0\_mem\_alloc} 调用成功。
\end{enumerate}

任一项不满足时，非 strict 作业会禁用 HotCache，状态访问仍由底层 SwissTable 正常服务；strict 作业则直接失败，保证要求 L0 的部署不会在缓存未激活时继续运行。运行时应结合日志与指标确认：\texttt{L0 cache disabled by config} 表示配置关闭；\texttt{L0 hardware not available ... forcibly disabled} 表示硬件门禁失败；\texttt{engine\_start active=1} 表示当前 TaskManager 首次成功创建进程级 manager；\texttt{engine\_reuse active=1} 表示同进程后续 StateEngine 正在复用该 manager；\texttt{state\_attach} 表示受支持状态已取得 HotSet（见 \S7.1）。

每个 TaskManager 进程只持有一个共享 \texttt{HotCacheManager}。因此 \texttt{l0-cache.total-size} 的预算划分按 TaskManager 进程数计算；提高 operator 并行度会增加共享池中的 state cache 数量，但不会再重复创建硬件 tuner。不同 TaskManager 之间仍是独立进程和独立 manager。

\subsection{SwissTable 容量参数}

以下参数决定单张 SwissTable 的隐含容量范围，默认值适合绝大多数作业，仅在超大规模或调优时才需要修改：

\begin{table}[H]
\centering
\small
\begin{tabularx}{\textwidth}{L{6.0cm} L{2.0cm} X}
\toprule
\textbf{配置项} & \textbf{默认值} & \textbf{说明} \\
\midrule
\cfg{state.backend.forl0.initial-table-capacity} & 64 & 初始 slot 数，必须为 2 的幂且 $\ge 8$。 \\
\cfg{state.backend.forl0.max-table-capacity}     & 0 & 单表容量上限；0 表示不设上限，非零值必须为 2 的幂。 \\
\cfg{state.backend.forl0.main-table.load-factor-threshold} & 0.875 & 触发 rehash / grow 的负载阈值，可在 0.5--0.875 范围内配置。 \\
\bottomrule
\end{tabularx}
\end{table}

\section{作业中的状态使用}

\project{} 使用 Flink 标准状态接口。用户作业中状态的声明、读写和清理方式与使用 HashMapStateBackend 时一致。

\subsection{DataStream API}

以下示例在 \texttt{KeyedProcessFunction} 中维护单值状态。切换到 \project{} 不需要对代码作任何修改：

\begin{lstlisting}[language=javalang]
StreamExecutionEnvironment env =
    StreamExecutionEnvironment.getExecutionEnvironment();

// 方式一：集群配置已将 state.backend.type 设为 org.apache.flink.state.forl0.ForL0StateBackendFactory，无需代码改动
// 方式二：在作业中显式设置
env.setStateBackend(new ForL0StateBackend());

input
    .keyBy(t -> t.f0)
    .process(new KeyedProcessFunction<String, Tuple2<String, Long>, String>() {
        private ValueState<Long> countState;

        @Override
        public void open(Configuration parameters) {
            countState = getRuntimeContext().getState(
                new ValueStateDescriptor<>("count", Long.class));
        }

        @Override
        public void processElement(Tuple2<String, Long> value,
                                   Context ctx,
                                   Collector<String> out) throws Exception {
            Long current = countState.value();
            current = (current == null) ? 1L : current + 1;
            countState.update(current);
            out.collect(value.f0 + ": " + current);
        }
    });
\end{lstlisting}

\subsection{Flink SQL 与 Table API}

SQL 作业通过集群配置启用 \project{} 即可，运行方式与使用其他状态后端时没有差别。例如：

\begin{lstlisting}[language=yamlcmd]
CREATE TABLE orders (
  order_id BIGINT,
  user_id  BIGINT,
  amount   DOUBLE,
  order_time TIMESTAMP(3),
  WATERMARK FOR order_time AS order_time - INTERVAL '5' SECOND
) WITH (...);

SELECT
  user_id,
  TUMBLE_START(order_time, INTERVAL '1' MINUTE) AS window_start,
  SUM(amount) AS total_amount
FROM orders
GROUP BY user_id, TUMBLE(order_time, INTERVAL '1' MINUTE);
\end{lstlisting}

SQL 作业会使用已配置的 \project{}。具体状态访问路径由运行时 serializer、key 和 namespace 类型决定，无需在 SQL 中增加 ForL0 专用语法。

\subsection{命名空间行为}

\project{} 的状态组织区分命名空间类型，对不同场景给出不同的访问路径：

\begin{itemize}
  \item \textbf{VoidNamespace 直通路径}：在非窗口算子中（例如 \texttt{KeyedProcessFunction} 中的 \texttt{ValueState}），命名空间恒为 \texttt{VoidNamespace}。\project{} 检测到这种情形后直接访问当前 KeyGroup 对应的 SwissTable，跳过命名空间容器，避免额外的容器查找开销。
  \item \textbf{TimeWindow 命名空间路径}：在窗口算子中，以 \texttt{TimeWindow} 作为命名空间的状态由每个命名空间独立的 SwissTable 承载。命名空间之间隔离，避免 key 冲突。
  \item \textbf{空命名空间回收}：当某个命名空间下的 SwissTable 因清理操作变为空时，\project{} 会将该命名空间条目从容器中移除，避免无效命名空间持续占用内存。
\end{itemize}

% ---------- 图：命名空间组织 ----------
\begin{figure}[htbp]
\centering
\begin{tikzpicture}[x=1cm,y=1cm]
  % VoidNamespace 路径
  \node[figtitle, anchor=west] at (0,4.2) {VoidNamespace 路径};
  \node[box, minimum width=2.6cm] (kg1) at (2.2,3.3) {KeyGroup[i]};
  \node[accentbox, minimum width=3.2cm] (tab1) at (6.1,3.3) {SwissTable\,<\,K, S\,>};
  \draw[line] (kg1.east) -- (tab1.west);

  % TimeWindow 路径
  \node[figtitle, anchor=west] at (0,2.0) {TimeWindow 命名空间路径};
  \node[box, minimum width=2.6cm] (kg2) at (2.2,1.1) {KeyGroup[i]};
  \node[box, minimum width=3.2cm] (nsmap) at (6.2,1.1) {Map<TimeWindow,\\SwissTable<K,S>>};
  \node[accentbox, minimum width=3.2cm] (tab2a) at (11.2,1.75) {SwissTable \#1};
  \node[accentbox, minimum width=3.2cm] (tab2b) at (11.2,0.45) {SwissTable \#2};
  \draw[line] (kg2.east) -- (nsmap.west);
  \draw[line] (nsmap.east) -- (tab2a.west);
  \draw[line] (nsmap.east) -- (tab2b.west);
\end{tikzpicture}
\caption{\project{} 的命名空间访问路径}
\label{fig:namespace}
\end{figure}

\section{Checkpoint 与 Savepoint}

\project{} 使用原生 Copy-On-Write 机制建立一致性视图，并通过 Flink Checkpoint 框架异步写出。用户侧的配置方式和操作方式保持不变。

\subsection{Checkpoint 配置}

\begin{lstlisting}[language=yamlcmd]
execution:
  checkpointing:
    interval: 60000           # 60 秒触发一次
    mode: EXACTLY_ONCE
    externalized-checkpoint-retention: RETAIN_ON_CANCELLATION

state:
  checkpoints:
    dir: file:///data/flink/checkpoints
  savepoints:
    dir: file:///data/flink/savepoints
\end{lstlisting}

\subsection{Savepoint 触发与恢复}

可通过 Flink CLI 触发 Savepoint 与恢复：

\begin{lstlisting}[language=shellcmd]
# 为运行中的作业触发 Savepoint
$FLINK_HOME/bin/flink savepoint <jobId> file:///data/flink/savepoints

# 从指定 Savepoint 启动作业
$FLINK_HOME/bin/flink run \
  -s file:///data/flink/savepoints/savepoint-xxx \
  -d your-job.jar
\end{lstlisting}

\subsection{兼容性约定}

\begin{itemize}
  \item 从 \project{} 产出的 Checkpoint 或 Savepoint，可在同一版本的 \project{} 之间互相恢复。
  \item 从其他状态后端迁移到 \project{} 时，应使用 canonical Savepoint，并在测试环境验证状态类型、serializer 与定时器恢复结果。其他后端生成的 Checkpoint 不作为跨后端迁移入口。
  \item 进行跨后端迁移前，建议在同等输入的测试环境下对账一次状态清单，以降低升级风险。
\end{itemize}

\section{作业提交}

\subsection{Flink CLI}

\begin{lstlisting}[language=shellcmd]
$FLINK_HOME/bin/flink run -d \
  -Dstate.backend.type=org.apache.flink.state.forl0.ForL0StateBackendFactory \
  -Dstate.backend.forl0.l0-cache.enabled=true \
  -p 8 \
  your-application.jar
\end{lstlisting}

如已在集群配置中将 \texttt{state.backend.type} 设为 \texttt{org.apache.flink.state.forl0.ForL0StateBackendFactory}，命令行的 \texttt{-Dstate.backend.type} 可省略。

\subsection{REST API}

\begin{lstlisting}[language=shellcmd]
# 上传 JAR
curl -X POST -H "Expect:" \
  -F "jarfile=@your-application.jar" \
  http://<jobmanager>:8081/jars/upload

# 提交作业并覆盖状态后端配置
curl -X POST http://<jobmanager>:8081/jars/<jar-id>/run \
  -d '{
        "programArgs": "--parallelism 8",
        "flinkConfiguration": {
          "state.backend.type": "org.apache.flink.state.forl0.ForL0StateBackendFactory"
        }
      }'
\end{lstlisting}

\section{运行监控}

\subsection{日志约定}

\project{} 相关的运行期日志一律以 \texttt{[ForL0]} 前缀输出。常见条目见下表：

\begin{table}[H]
\centering
\small
\begin{tabularx}{\textwidth}{L{7.4cm} X}
\toprule
\textbf{日志} & \textbf{含义} \\
\midrule
\texttt{[ForL0] Native engine loaded via java.library.path} & 通过 \texttt{java.library.path} 加载成功 \\
\texttt{[ForL0] Native engine loaded from bundled resource: \ldots} & 从 JAR 内置资源解压并加载成功 \\
\texttt{[ForL0] Creating ForL0StateBackend from configuration.} & 通过 Factory 从配置创建后端 \\
\texttt{[ForL0] Building ForL0KeyedStateBackend with C++ engine\ldots} & KeyedStateBackend 构造开始 \\
\texttt{[ForL0] C++ engine created: handle=\ldots, l0Enabled=\ldots} & C++ 引擎创建完成，l0Enabled 反映 HotCache 是否激活 \\
\texttt{[ForL0-HotCache] L0 cache disabled by config.} & 当前配置未启用 HotCache \\
\texttt{[ForL0-HotCache] engine\_start active=1 \ldots} & 当前 TaskManager 首次成功初始化进程级 L0 manager \\
\texttt{[ForL0-HotCache] engine\_reuse active=1 \ldots} & 同一 TaskManager 内其他 StateEngine 复用已有 manager，属于正常并行行为 \\
\texttt{[ForL0-HotCache] state\_attach name=\ldots} & 受支持状态已从共享池取得 HotSet \\
\texttt{[ForL0-HotCache] WARN: L0 hardware not available \ldots} & 配置要求 L0，但设备、库或分配门禁失败 \\
\texttt{[ForL0] ForL0KeyedStateBackend built successfully.} & KeyedStateBackend 构造完成 \\
\texttt{[ForL0] Failed to load native engine: \ldots} & 原生库加载失败，后端不会启动 \\
\bottomrule
\end{tabularx}
\end{table}

\subsection{运行期加载检查 API}

在用户代码或自检脚本中判断原生引擎是否已加载：

\begin{lstlisting}[language=javalang]
import org.apache.flink.state.forl0.NativeEngine;

boolean loaded = NativeEngine.isLoaded();
NativeEngine.ensureLoaded();  // 未加载时抛出 RuntimeException
\end{lstlisting}

原生引擎本身在所在类被容器加载时完成一次性的 \texttt{System.loadLibrary("forl0\_engine")}，不提供重新初始化接口。

\subsection{HotCache 指标}

启用 L0 Hot-Key Cache 且设置 \texttt{forL0.metricsCollector.enabled=true} 时，后端在作业的 \texttt{MetricGroup} 下注册以下 gauge，指标名为 \texttt{<operator>.forl0.hotcache.<field>}，可直接在 Flink Web UI 或对接的指标后端（Prometheus / InfluxDB 等）中观察：

\begin{table}[H]
\centering
\small
\begin{tabularx}{\textwidth}{L{3.8cm} X}
\toprule
\textbf{字段} & \textbf{含义} \\
\midrule
\texttt{active}         & HotCacheManager 是否真正激活（1 表示硬件门禁通过；0 表示未激活） \\
\texttt{bytesCapacity} & L0 分区从 \texttt{l0\_mem\_alloc} 实际取得的总容量，单位为字节 \\
\texttt{bytesUsed}     & 当前 HotSet 已用字节数 \\
\texttt{totalSets}     & 全局 HotSet 总数 \\
\texttt{freeSets}      & 当前可用的 HotSet 数 \\
\texttt{lookups}       & HotCache 查询总次数 \\
\texttt{hits}          & HotCache 命中总次数（命中率由 \texttt{hits / lookups} 给出） \\
\texttt{invalidations} & HotCache 失效总次数（erase / 覆写触发） \\
\texttt{bytesRequested} & 当前 TaskManager 请求的 L0 配额，单位为字节 \\
\texttt{writes} & HotCache 写入总次数 \\
\texttt{writeBypassEvents} & 连续写入触发缓存准入暂停的次数 \\
\texttt{stateCaches} & 当前已附着的状态缓存数量 \\
\bottomrule
\end{tabularx}
\end{table}

当设置 \texttt{l0-cache.enabled=true} 但 \texttt{active} 始终为 0 时，说明硬件门禁未通过（见 \S3.3），非 strict 作业已回退到纯 SwissTable 路径；此时 \texttt{bytesCapacity=0}、\texttt{hits=0}、\texttt{lookups=0}。若 \texttt{active=1} 且存在 \texttt{state\_attach}，但 \texttt{lookups=0}，表示缓存已附着但当前负载没有执行缓存读取，命中率应记为“不适用”而不是 0\%。

\section{调优建议}

\subsection{并行度与 KeyGroup}

\project{} 依照 Flink 标准的 KeyGroup 路由模型组织状态。建议并行度设置与实际可用 slot 数对齐，避免因并行度过高导致的 KeyGroup 过细、快照开销相对增大的问题。鲲鹏服务器常见的推荐公式为：

\[
\text{parallelism} = \text{TaskManager 数量} \times \text{每 TM 的 slot 数}
\]

\subsection{L0 Cache 容量}

开启 L0 Cache 的作业可参考以下容量范围。实际分配成功的字节数由 \texttt{forl0.hotcache.bytesCapacity} 反映：

\begin{table}[H]
\centering
\small
\begin{tabularx}{\textwidth}{L{3.4cm} L{4.2cm} X}
\toprule
\textbf{配置场景} & \textbf{建议配置} & \textbf{说明} \\
\midrule
单 TaskManager & \cfg{l0-cache.size}: 20\,mb & 使用兼容配置项设置当前进程配额。 \\
多 TaskManager & \cfg{l0-cache.total-size} + \cfg{expected-engines} & 按设备可用预算设置总量，并按 TaskManager 进程数均分。 \\
单状态缓存 & \cfg{l0-cache.state-size}: 1\,mb & 仅符合类型条件的标量 ValueState 会申请该配额。 \\
\bottomrule
\end{tabularx}
\end{table}

\subsection{TaskManager 内存}

\project{} 的 Keyed KV 状态位于原生内存，优先队列仍由 Flink 管理；本后端不使用 Flink Managed Memory 保存 KV 状态。据此可作如下调整：

\begin{itemize}
  \item 适当增加 \texttt{taskmanager.memory.process.size} 的上限，以容纳原生状态数据。
  \item 相应降低 \texttt{taskmanager.memory.managed.size}，为原生数据面让出空间。
  \item 若开启 L0 Cache，确保 \texttt{l0-cache.total-size} 与设备实际可用预算匹配，并正确设置 TaskManager 进程数。
\end{itemize}
\subsection{Checkpoint 间隔}

\project{} 通过 Copy-On-Write 缩短同步准备阶段，快照开销主要来自变更记录、序列化与写出。建议：

\begin{itemize}
  \item 保持 \texttt{state.backend.forl0.async-snapshots: true}。
  \item 根据容灾目标设置 \texttt{execution.checkpointing.interval}，无特殊需求时使用 30--60 秒量级。
  \item 生产环境将 Checkpoint 目录指向稳定的共享存储（HDFS、NAS 等）。
\end{itemize}

\section{常见问题}

\begin{table}[H]
\centering
\small
\begin{tabularx}{\textwidth}{L{4.8cm} L{3.4cm} X}
\toprule
\textbf{现象} & \textbf{可能原因} & \textbf{处理方式} \\
\midrule
\texttt{UnsatisfiedLinkError} & 原生库路径未配置 & 检查 \texttt{java.library.path} 是否包含 \texttt{libforl0\_engine.so} 所在目录 \\
\texttt{ClassNotFoundException} & JAR 未部署 & 确认 \texttt{\$FLINK\_HOME/lib/} 下存在 \cfg{flink-statebackend-forL0-*.jar} \\
日志中 \texttt{l0Enabled=false} & L0 硬件/库不可用或未启用 & 核查 \texttt{/dev/hisi\_l0} 或 \texttt{/dev/l0}、\texttt{libl0mempool.so} 以及 \texttt{l0-cache.enabled} \\
Checkpoint 失败 & Checkpoint 目录权限不足 & 确认 \texttt{state.checkpoints.dir} 路径对 Flink 进程可写 \\
\cfg{forl0.hotcache.active=0} 但配置已开启 & 硬件门禁未通过 & 检查 L0 设备节点是否存在，\texttt{libl0mempool.so} 是否在搜索路径上 \\
配置要求 L0，但日志出现 \texttt{disabled by config} & 配置未生效或被作业参数覆盖 & 检查集群配置与作业提交参数中的 \texttt{l0-cache.enabled} \\
只有 \texttt{engine\_start} 后又有多个 \texttt{engine\_reuse} & 同一 TaskManager 内并行 StateEngine 共享 manager & 正常行为；继续确认存在 \texttt{state\_attach}，且无 hardware-not-available 警告 \\
作业 OOM & 状态规模超出进程内存预算 & 增加 \cfg{taskmanager.memory.process.size}，或设置状态 TTL \\
\bottomrule
\end{tabularx}
\end{table}

\end{document}
